\documentclass[12pt]{article}

\usepackage[margin=1in]{geometry}
\usepackage{times}
\usepackage[T1]{fontenc}
\usepackage{parskip}

\pagestyle{empty}

\begin{document}

\textbf{Mentoring Plan}

This project will support one graduate research assistant (GRA) for five years. Mentorship of the GRA will be done by PI Poulsen. Poulsen has thus far mentored 3 PhD students and 5 MS students, having funded two of them with his current NSF project (2434364).
% Collectively, his students have co-authored 2 peer-reviewed journal article and led 3 conference proceedings, with many more papers in preparation.

Part of the strength of Poulsen's mentorship lies in his focus on building community. He has cultivated a supportive community through the computer science education research group at Utah State University, which he co-leads with two other faculty in the department. The computer science education research group meets weekly to discuss research ideas and for students to get feedback on their work. Thus, the mentoring environment for the GRA will be very strong. Specific elements of the mentoring plan include:

1. In the first week, Poulsen will meet with the GRA to create an individualized development plan. They will co-identify development activities and milestones, such as leading manuscripts and papers for conferences where they would like visibility. This plan will be reviewed and updated each year of the project. The GRA and PI will craft a work plan for the GRA at the start of each semester. The GRA and PI will discuss progress at their weekly one-on-one meeting as well as strategies for meeting the next week's milestones.

2. The GRA will be encouraged to attend seminars and workshops for professional development. The Office of Research at USU has a research support webpage and mailing list with information and resources for graduate students. They offer workshops on mentoring undergraduate researchers, research presentations, funding opportunities, teaching, and grant writing. The GRA will be encouraged to attend virtual and conference workshops such as those offered through AERA Professional Development and Training.

3. Poulsen will mentor his GRA on data collection, statistical data analysis, qualitative analysis and writing. The GRA will be mentored in leading publications in peer-reviewed journals and presenting at national and international conferences. Funds for the GRA's travel to conferences are included in the budget request for this project. The PI will also mentor the GRA on the grant writing process.

4. The PI and GRA will have weekly meetings for the GRA to check in, provide updates, and receive direct, formative feedback. This weekly meeting will be in addition to the weekly project team meetings and the PI's lab meetings where current research is discussed.

5. The PI will encourage the GRA to develop a network of mentors and help them to identify and connect with other scholars with expertise relevant to their areas of interest. Success of this mentoring plan will be evaluated by tracking the progress of the GRA according to their individualized development plan, through annual performance reviews, and through an exit interview at the completion of the appointment. The GRA subsequent activities and career trajectory will be included in the final project report.

\end{document}
